% SLICE-UNSAT REPORT::HYPERGEOM-MOMENT1::2026-10-04T22:13-07:00
\documentclass{article}
\usepackage{amsmath,amssymb,amsthm,hyperref}
\newtheorem{lemma}{Lemma}
\begin{document}
\section*{HYPERGEOM-MOMENT1}
\textbf{Status: OPEN.} The requested Git ref INTERMEDIATE1 is not a branch in the repository, so this run used the newest default-branch state and the INTERMEDIATE1 state recorded in TASK-008. AGENTS.md, CODEX.md, the pnp-research skill, root main.tex, and TASK-008 were read first.

\paragraph{Literature.}
Primary sources checked before the claim below: Cavalar et al., \emph{Derandomized Sunflowers and Radical Lower Bounds}, ECCC TR26-220 (2026), \url{https://eccc.weizmann.ac.il/report/2026/220/}; Cavalar--Kumar--Rossman, \emph{Monotone Circuit Lower Bounds from Robust Sunflowers}, Algorithmica 84 (2022), DOI 10.1007/s00453-022-01000-3; Amireddy et al., \emph{Low-Degree Testing Over Boolean Slices}, ECCC TR26-152 (2026), \url{https://eccc.weizmann.ac.il/report/2026/152/}.

\begin{lemma}[DISJOINT-SLICE-RATIO]
\textbf{PROVED/EXACT.}
Let $X$ be uniform on $\{x\in\{0,1\}^M:|x|=t\}$. For disjoint coordinate sets $A,B$ of sizes $a,b$, with $a+b\le t$,
\[
R_{a,b}:=\frac{\Pr[X_A=1,X_B=1]}{\Pr[X_A=1]\Pr[X_B=1]}
=\prod_{j=0}^{b-1}\frac{1-a/(t-j)}{1-a/(M-j)}\le1.
\]
If $a+b<t+1$, then
\[
0\le-\log R_{a,b}\le\frac{ab}{t-a-b+1},
\qquad
0\le1-R_{a,b}\le\frac{ab}{t-a-b+1}.
\]
\end{lemma}
\paragraph{Proof.}
Sampling without replacement gives
$\Pr[X_A=1]=(t)_a/(M)_a$ and
$\Pr[X_B=1\mid X_A=1]=(t-a)_b/(M-a)_b$, yielding the product. Each factor is at most one. Then
\[
-\log R_{a,b}\le\sum_{j=0}^{b-1}-\log(1-a/(t-j))
\le\sum_{j=0}^{b-1}\frac{a}{t-j-a}
\le\frac{ab}{t-a-b+1}.
\]
Use $1-R\le-\log R$.

\paragraph{SAT-SUNFLOWER2 consequence.}
At $M=8\binom q3$ and $t=\lfloor3.5q\rfloor$, disjoint exposed blocks of sizes $a,b$ have relative factorization error $O(ab/t)$ whenever $a+b=o(t)$. This exactly identifies the correction required where a product-measure moment proof factors disconnected components.

\paragraph{HYPERGEOM-PARTITION.}
\textbf{CANDIDATE, not proved.}
A width-$w$, $p$th-moment expansion exposes at most $r=pw$ coordinates. Sequential application suggests a total distortion $\exp(O(r^2/t))$, hence the candidate window $pw=o(\sqrt t)$. The missing step is the centered-moment partition sum in TR26-220. Centering can amplify error, so the conjunction ratio above is insufficient by itself.

\paragraph{Hostile controls.}
The lemma is generic to every fixed-weight slice and therefore also holds for XOR-SLICE-EASY. It cannot itself imply SAT hardness. LOCALPATTERN-GUARD, the P/poly positive-family barrier, and SLICELIFT2 remain active. The withdrawn switching-lemma route remains \textbf{FALSIFIED/RETIRED}. MOMENTBALANCE3 and ORDER3EXACT1 remain finite \textbf{COMPUTATIONAL} evidence only. MATCHED-MUF2 remains \textbf{OPEN}.

\paragraph{Next target.}
Transcribe the exact centered-moment partition identity from TR26-220 and either prove an overall $\exp(O((pw)^2/t))$ hypergeometric distortion or find the smallest partition falsifying that scale. No restricted-model result is promoted to unrestricted SAT runtime. Since load-bearing circuit approximation and positive-family high-order control remain \textbf{OPEN}, P versus NP remains open.

% END_OF_VERIFIED_REPORT::SLICE-UNSAT::HYPERGEOM-MOMENT1
\end{document}
